\documentclass[11pt]{article}
\usepackage[margin=1in]{geometry}
\usepackage{amsmath,amssymb,amsthm,hyperref}
\newtheorem{theorem}{Theorem}\newtheorem{lemma}[theorem]{Lemma}
\newtheorem{corollary}[theorem]{Corollary}
\newcommand{\E}{\mathbb E}\newcommand{\R}{\mathbb R}
\newcommand{\one}{\mathbf1}
\title{A logarithmic separation between local smoothing and rare neural features}
\author{ChatGPT-6 Astra\thanks{Independent AI research. Sole author.}}
\date{21 November 2026; working proof draft}
\hypersetup{hidelinks,pdftitle={A logarithmic separation between local smoothing and rare neural features},pdfauthor={ChatGPT-6 Astra}}
\emergencystretch=3em
\begin{document}\maketitle
\begin{abstract}
In a full-label Gaussian neural posterior, softplus smoothing changes the
rarity of additional macroscopic neurons before it changes their local
fluctuation law to the ReLU law. We prove a joint width--temperature limit:
when the inverse temperature is of order $\sqrt N/\log N$, the posterior
probability of each fixed extra-neuron count has a continuously shifted
polynomial exponent. The shift is an explicit square-root splitting cost.
The local critical-neuron marks at this scale remain Gaussian, with
transverse size $\sqrt{\log N/N}$; their ReLU crossover instead occurs
at inverse temperature of order $\sqrt N$. In six input dimensions, we
also identify an inverse-temperature covariance correction of rank 38,
caused by the critical large-neuron tail; in even dimensions at least eight,
the covariance instead has a regular second-order response. A boundary integral identifies
the regular second-order normalizer change, which produces an explicit
negative logarithm-squared evidence correction in the joint limit.
A positive Fourier integral shows more generally that removing softplus
smoothing increases every even moment of the Gaussian neuron's parity
correlation and strictly increases its critical normalizer. For the
input-centered random-feature and initialization tangent kernels,
sharpening also shifts normalized variance toward higher Boolean Fourier
degrees, strictly reducing the normalized RKHS cost of full parity.
The proof combines a uniform
corner estimate, a separate fine cutoff for bulk localization, and the
count-and-rank volume calculation. These are fixed-dimension, full-cube
posterior results, not a growing-dimension learning or sampling guarantee.
\end{abstract}

\section{Model and joint limit}
Fix even $r\ge6$, $X=\{-1,1\}^r$, $M=2^r$, $K=M/2$, and
$\chi(x)=\prod_i x_i$. Write $\bar x=(1,x)$ and
$\phi_\beta(t)=\beta^{-1}\log(1+e^{\beta t})$.
A neuron $\theta=(a,U)\in\R^{r+2}$ has standard Gaussian law $P_0$ and
output $\Phi_{\beta,x}=a\phi_\beta(U\cdot\bar x)$.
The width-$N$ output is $F_N=N^{-1}\sum_i\Phi_{\beta_N,i}$.
Set
\[
 \rho=\frac{\binom{r-2}{(r-2)/2}}{2^{r-1}\sqrt r},\quad
 p_0=\frac{\binom r{r/2}}{2^r},\quad
 \kappa_0=\frac{p_0\log2}{\rho},\quad
 Z_\beta=M^{-1}\chi^\top\Phi_\beta,
 \qquad Q_\beta=\mathcal M_\beta^{-1}e^{Z_\beta/\rho}P_0.
\]
Let $m_\beta=\E_{Q_\beta}Z_\beta$. Subscript $\infty$ denotes the ReLU
critical law and its normalizer and moments. The earlier critical-law
companions establish their finiteness and the critical-ray geometry.
The uniformity in temperature required here is proved below.
Assume
\[
        a_N=\frac{\sqrt N}{\beta_N\log N}\longrightarrow
                       a\in[0,\infty).                 \tag{1}
\]
In particular $\beta_N\to\infty$. Condition either on the hard event
$H_N=\{\chi(x)F_N(x)>s\text{ for all }x\}$, with $s>m_\infty$,
or use a likelihood proportional to
$\exp\{-N\lambda M^{-1}\sum_x\ell(\chi(x)F_N(x))\}$.
Here $\ell(y)=(1-y)^2/2$ or $\log(1+e^{-y})$, and in the soft case
$s$ is the minimizer of $g(y)=y/\rho+\lambda\ell(y)$, assumed larger
than $m_\infty$. Denote the posterior by $\Pi_N$ and put
\[
 \alpha=\frac{s-m_\infty}{K\rho},\quad
 \nu=\frac{r+1}{2},\quad b=\nu-2=\frac{r-3}{2},\quad
 C_H=\frac{Kr}{2},\quad C_L=\frac{K(r-2)}{2}.
\]
These exponents use the ReLU corner power, not the fixed-temperature
softplus power. Let $L_N(d)=\#\{i:\|\theta_i\|^2>dN\}$.

\begin{theorem}[Polynomial exponents at logarithmic smoothing]\label{thm:main}
Under (1), the hard evidence satisfies
\[
 P_0^N(H_N)=
 e^{-N(s/\rho-\log\mathcal M_{\beta_N})}
 N^{-C_H-a\kappa_0K\sqrt\alpha+o(1)}.                  \tag{2}
\]
The soft evidence, namely the prior expectation of the likelihood, is
\[
 e^{-N(g(s)-\log\mathcal M_{\beta_N})}
 N^{-C_L-a\kappa_0K\sqrt\alpha+o(1)}.                  \tag{3}
\]
Fix $J\ge1$ and $d>0$ with $\delta=d/2<\alpha/(J+2)$. For $1\le j\le J$,
\begin{align}
 \Pi_N\{L_N(d)\ge K+j\}
   &=N^{-jb-a\kappa_0\Delta_j+o(1)},\notag\\
 \Pi_N\{L_N(d)=K+j\}
   &=N^{-jb-a\kappa_0\Delta_j+o(1)},\qquad
 \Delta_j=j\sqrt\delta+\sqrt{\alpha-j\delta}-\sqrt\alpha>0.
                                                               \tag{4}
\end{align}
The dimension, conditioning strength, $J$, and $d$ are fixed. Equations
(2)--(4) identify logarithmic exponents, not asymptotic prefactors.
\end{theorem}

The exact normalizer $\mathcal M_{\beta_N}$ must be retained in
(2)--(3). Its convergence to $\mathcal M_\infty$ alone does not justify
replacing it inside a term multiplied by $N$. When $a=0$, (4) recovers
the ReLU logarithmic count exponents. This does not assert convergence
of the whole posterior to the ReLU posterior.

\section{A corner estimate uniform in inverse temperature}
For $z\in X$ define $v_z=(0,z/\sqrt r)$,
$\sigma_z=(-1)^{r/2+1}\chi(z)$, and
$f_z(x)=\sigma_z\max\{z\cdot x/\sqrt r,0\}$.
In fixed small disjoint cones about $(\sigma_z,v_z)$ use coordinates
\[
 a=\sigma_z A,\quad U=Rv_z+t/\sqrt c,\quad
 c=AR>0,\quad v=A-R,\quad t\perp v_z.
\]
The cone keeps $q=A/\sqrt c$ between positive fixed constants.
Let $q_{\beta,z}(c)$ be the $Q_\beta$ density of $c$ in this cone,
with the other coordinates integrated out.

\begin{lemma}[Uniform radial estimate]\label{lem:corner}
There are constants $c_0,C_1,C_2>0$, independent of $\beta\ge1$, such
that, for $c\ge c_0$,
\[
 C_1 c^{-\nu}\left[1+\left(\frac{\sqrt c}{\beta}\right)^{r/2}\right]
 e^{-\kappa_0\sqrt c/\beta}
 \le q_{\beta,z}(c)\le
 C_2 c^{-\nu}\left[1+\left(\frac{\sqrt c}{\beta}\right)^{r/2}\right]
 e^{-\kappa_0\sqrt c/\beta}.                          \tag{5}
\]
All ray cones have the same constants. The critical parameter laws have
uniform tails
\[
 \sup_{\beta\ge1}Q_\beta\{\|\theta\|^2>u\}
                       \le Cu^{-(\nu-1)},\qquad u\ge1. \tag{6}
\]
Their normalizers, output means, and output covariances converge to the
ReLU values as $\beta\to\infty$.
\end{lemma}
\begin{proof}
Write $\zeta=\beta/\sqrt c$ and
\[
 V_\zeta(t)=\frac1{\zeta\rho M}
       \sum_{x:z\cdot x=0}\log\left(2\cosh\frac{\zeta t\cdot\bar x}{2}\right),
 \qquad H_\zeta(t)=V_\zeta(t)-\frac{\kappa_0}{\zeta}.
 \tag{7}
\]
The exact corner calculation gives, for the negative log of the
unnormalized Gaussian tilted density,
\[
 D_\beta=\frac{v^2}{2}+\frac{\|t\|^2}{2c}
                    +qV_\zeta(t)+O(A\beta^{-1}e^{-c_1\beta R}). \tag{8}
\]
The error is uniformly bounded and tends to zero uniformly for large
$c$ and $\beta\ge1$. The zero-layer vectors span the transverse space.
The scalar bounds for $\log\cosh u$ at zero and infinity therefore give
\[
 c_2\min\{\zeta\|t\|^2,\|t\|\}
 \le H_\zeta(t)\le
 C_3\min\{\zeta\|t\|^2,\|t\|\}.                    \tag{9}
\]
Since $A\ge\sqrt c+v/2$,
\[
       \frac{v^2}{2}+q\frac{\kappa_0}{\zeta}
       \ge\frac{v^2}{4}+\frac{\kappa_0}{\zeta}
                                      -\frac{C}{\beta^2}.
\]
The Jacobian is $c^{-r/2}/\sqrt{v^2+4c}\le c^{-\nu}/2$.
Integration of (9) yields at most $C(1+\zeta^{-r/2})$.
For the lower bound restrict $v$ to a fixed interval and $t$ to a
fixed ball when $\zeta\ge1$, or to a ball of radius
$c_3\zeta^{-1/2}$ when $\zeta<1$. These sets lie in the cone for
large $c$, and their exponent beyond $\kappa_0/\zeta$ is uniformly
bounded. Their transverse volume is comparable to
$1+\zeta^{-r/2}$.

For clarity, normalizer bounds do not require the conclusion being
proved. On these cones, $\beta^{-1}\log(2\cosh(\beta u/2))\ge|u|/2$
shows that the soft critical density is at most a constant times the
ReLU critical density. Off the cones the ReLU exponent has a fixed
positive quadratic gap. The uniform approximation
$|\phi_\beta(t)-\max(t,0)|\le\log2/\beta$ changes that exponent by at
most $C\|\theta\|/\beta$, which is absorbed by half the gap and a
constant. On a bounded parameter region the density is uniformly
bounded. Thus all unnormalized densities have a common integrable
majorant, while $\mathcal M_\beta\ge1$ by integrating the Gaussian
readout. This supplies uniform upper and lower normalizer bounds.

In (5), $(1+w^{r/2})e^{-\kappa_0 w}$ is bounded for $w\ge0$.
Integrating $c^{-\nu}$ on the cones and adding the off-cone Gaussian
tail proves (6). It also gives uniform parameter moments of each order
strictly below $r-1$. In particular, a moment of order $4+\epsilon$
exists for some $\epsilon>0$, since $r\ge6$. Pointwise convergence of
the densities, uniform integrability, and
$\|\Phi_\beta\|\le C(1+\|\theta\|^2)$ prove convergence of the
normalizers and the first two output moments.
\end{proof}

Along (1), for $\eta$ in a fixed compact subset of $(0,\infty)$,
(5) implies uniformly
\[
 q_{\beta_N,z}(N\eta)
 =N^{-\nu}\exp\{-\kappa_0\sqrt{N\eta}/\beta_N\}\,N^{o(1)}
 =N^{-\nu-a\kappa_0\sqrt\eta+o(1)}.                 \tag{10}
\]
The extra mark volume is at most a fixed power of $\log N$.
On such macro amplitudes, (8)--(9) also show that marks outside
$|v|^2+\|t\|\le B\log N$ have arbitrarily small polynomial weighted
mass by increasing the fixed $B$. On the retained marks,
\[
        \Phi_{\beta_N}=cf_z+O(\log N),\qquad
        \|\theta\|^2=2c+O(\log N).                  \tag{11}
\]
The smoothing contribution to the first error is bounded by
$CA/\beta_N=O(\log N)$.

\section{Uniform bulk estimates and a separate fine cutoff}
There is a fixed convolution block of the $Q_\beta$ output laws, for
all sufficiently large $\beta$, that minorizes the same centered
isotropic Gaussian law with a common positive mass. To see this, choose
bounded incoming neighborhoods whose ReLU output columns isolate the
cube points. They remain full rank for large $\beta$. Conditional
readouts are independent Gaussians with means $G_\beta(U)/\rho$ and
unit variance. On those neighborhoods their output covariances are
uniformly positive definite and bounded, and their means are bounded.
Their densities consequently dominate a smaller centered isotropic
Gaussian density. The incoming-neighborhood probabilities are bounded
below by the uniform critical-density estimates.

Independent block mixtures now give two uniform consequences. Apart
from a subprobability of mass $Ce^{-c n}$, a sum of $n$ outputs is a
mixture of translates of isotropic Gaussians with covariance between
$cnI$ and $CnI$. Moreover, for fixed $A,h>0$,
\[
 \inf_{\|w\|\le A\sqrt n}
 Q_{\beta_N}^{n}\left\{
 \left\|\sum_i\Phi_i-nm_{\beta_N}\chi-w\right\|_\infty<h
 \right\}\ge c_{A,h}n^{-M/2}.                       \tag{12}
\]
Uniform second moments keep the mixture centers within $O(\sqrt n)$
with positive probability; the Gaussian component then gives (12).
These statements use no quantitative central limit theorem.

Changing the prior to $Q_{\beta_N}^N$ leaves a residual $0\le W_N\le1$.
With $Y=\sum_i\Phi_i$, it is
\[
 e^{-\sum_x(\chi(x)Y_x-Ns)/(M\rho)}\prod_x\one_{\{\chi(x)Y_x>Ns\}}
 \quad\hbox{or}\quad
 e^{-N M^{-1}\sum_x[g(\chi(x)Y_x/N)-g(s)]}.            \tag{13}
\]
Its weighted mass outside $\|Y-Ns\chi\|\le\eta N$ is exponentially
small for each fixed $\eta>0$.

The target cutoff $d$ in (4) stays fixed. Introduce a different,
smaller cutoff $eN$ solely for the proof. Uniform tails and variances,
and Bennett's inequality applied coordinatewise to
$\Phi_i\one_{\{\|\theta_i\|^2\le eN\}}$, show that for every
$\eta,A_0>0$ one can choose fixed $e>0$ sufficiently small so that
\[
 Q_{\beta_N}^N\left\{
 \left\|\sum_i\Phi_i\one_{\{\|\theta_i\|^2\le eN\}}
            -Nm_{\beta_N}\chi\right\|>\eta N\right\}
                                         \le C N^{-A_0}. \tag{14}
\]
Indeed the truncation bias is $o(N)$ uniformly, the coordinate bound
is $CeN$, and the total variance is $O(N)$. The Bennett exponent is
at most $-(c\eta/e)\log N+O_{e,\eta}(1)$.
Also
\[
             Q_{\beta_N}^N\{L_N(e)>D\}
                             \le C_{e,D}N^{-(D+1)b}.    \tag{15}
\]
Thus $D$ may be chosen large enough to make this error arbitrarily
small polynomially. This separation of cutoffs matters: driving the
localization tolerance to zero at the fixed target cutoff $d$ would
not justify an arbitrarily accurate Bennett bound.

On the resulting coarse event, (11), (13), and (14) imply that every
fine-large neuron is in a critical cone, with amplitude between fixed
positive multiples of $N$, and that its pair totals satisfy
\[
       \sum_{j:z_j\in\{z,-z\}}\frac{c_j}{N}
                            =\alpha+O(\eta)+o(1).       \tag{16}
\]
Here the even Walsh block of the ray columns is invertible and sends
$(s-m_\infty)\chi$ to $\alpha\one$; parity also bounds the sum of all
amplitudes. Every antipodal pair occurs. Off-cone errors are exponentially
small. The mark cutoff in (11) makes the remaining errors arbitrarily
small polynomially. The number of fine-large neurons is between $K$
and $D$, and each associated output matrix has rank at least $K$.

\section{Count, rank, and square-root cost}
For a specified list of $L$ fine-large rays let $A$ be its output matrix
and $R_A$ its rank. Coarea in the amplitude box, the Gaussian component
of the ordinary sum, and the at most $N^L$ choices of labels give the
same volume power as for ReLU:
\[
 N^{-L\nu}\,N^{L-R_A}\,N^{-(M-R_A)/2}\,N^\xi\,N^L
       =N^{-Lb-(M+R_A)/2+\xi},\qquad
 \xi=\begin{cases}0&\text{hard},\\M/2&\text{soft}.\end{cases} \tag{17}
\]
The fiber in the macro box has volume $O(N^{L-R_A})$. Integration of
the ordinary Gaussian density along the range of $A$ leaves a marginal
bounded by $CN^{-(M-R_A)/2}$. The integral of the residual in (13) is
$O(N^\xi)$. Replacing it by its supremum over the $O(\log N)$ error
box from (11) changes this only by a power of $\log N$. The exponentially
small bad Gaussian component is negligible. The mark integrations in
(10) likewise contribute only $N^{o(1)}$.

The additional factor from the radial density is
\[
       N^{-a\kappa_0\sum_{i=1}^L\sqrt{c_i/N}+o(1)}.    \tag{18}
\]
At exact pair totals $\alpha$, its sum is at least $K\sqrt\alpha$.
If at least $K+j$ amplitudes exceed $\delta$, it is at least
$K\sqrt\alpha+\Delta_j$. Here is the elementary optimization, including
possible uncounted pairs. If a pair has $n\ge1$ counted pieces, merging
all other pieces into the largest one and moving mass to the boundary
shows that its least square-root sum is
$(n-1)\sqrt\delta+\sqrt{\alpha-(n-1)\delta}$.
A pair with no counted piece costs at least $\sqrt\alpha$.
If $h$ pairs are uncounted, there are at least $j+h$ excess counted
pieces among the occupied pairs. Retain any $j$ of these excess pieces.
The increments $\Delta_l$ are increasing and concave on $0\le l\le j$,
and $\Delta_0=0$, so splitting those $j$ pieces between pairs cannot
reduce their total cost below $\Delta_j$.
For (16) the same lower bound has an error tending to zero as
$\eta\downarrow0$, by continuity on the finite compact amplitude boxes.
The strict norm cutoff agrees with the amplitude cutoff up to $o(1)$
by (11).

Since $L\ge K$, $R_A\ge K$, (17)--(18) prove the upper logarithmic
evidence bound. Under the extra-count event $L\ge K+j$, they prove
an upper exponent $C_H+jb+a\kappa_0(K\sqrt\alpha+\Delta_j)$ or its
soft analogue. Choose the remainder order $A_0$ larger than these
fixed exponents before choosing $e,D,B$. There are finitely many fine
ray lists. First let $N\to\infty$, then $\eta\downarrow0$.

For the matching lower bound choose a balanced set of $K$ rays, with
one ray per pair. For evidence alone, give each amplitude a $\sqrt N$
interval about $[Ns-(N-K)m_{\beta_N}]/(K\rho)$ and restrict marks to
fixed bounded boxes. For the count event, duplicate one ray $j$ times.
Its $j$ small amplitudes are just above $N\delta$, the remaining
amplitude is near $N(\alpha-j\delta)$, and all other pair totals are
unchanged. The amplitude region has $K$ directions of length
$\sqrt N$ and $j$ free directions of length proportional to $N$.
Choose its fixed split neighborhood so that the square-root sum is
within an arbitrarily small constant of $K\sqrt\alpha+\Delta_j$.
This is possible because $\delta<\alpha/(j+1)$.

The selected output differs from the desired mean by $O(\sqrt N)$;
the smoothing error $O(\log N)$ is smaller. Equation (12) gives the
hard lower small-box factor $N^{-M/2}$ with residual bounded below.
For soft loss, uniform second moments give a positive probability
of bulk fluctuations $O(\sqrt N)$, again with residual bounded below.
Summing the order $N^{K+j}$ label assignments yields the matching
powers in (17). On $L_N(e)\le D$ their overlap is bounded by a constant;
the discarded sum is at most $N^{K+j}Q_{\beta_N}^N\{L_N(e)>D\}$,
which is negligible after increasing $D$. Let the split neighborhood
approach its boundary to obtain the claimed logarithmic lower bound.
The case $j=0$ gives the lower evidence bound. Restoring the prior
factor proves (2)--(3), and dividing gives the first formula in (4).
The same tail argument for $j+1$ has strictly larger exponent, since
$b>0$ and $\Delta_{j+1}>\Delta_j$. Subtraction proves the exact-count
formula. The hypothesis $\delta<\alpha/(J+2)$ permits that final tail.

\section{Local marks at the logarithmic threshold}
Define, on $v_z^\perp$,
\[
                 Q_z=M^{-1}\sum_{x:z\cdot x=0}\bar x\bar x^\top.
\]
Its eigenvalues are $p_0$ in the bias direction and $p_0r/(r-1)$ in
the other $r-1$ directions.

\begin{theorem}[One-neuron Gaussian marks]\label{thm:marks}
Suppose $a\in(0,\infty)$ in (1), and fix $\eta>0$. Conditional on
$c=N\eta$ and a critical cone, write the physical transverse parameter
as $u=\sqrt{\log N/N}\,t$. Then $(v,t)$ converges in total variation to
independent Gaussians
\[
            v\sim N(0,1),\qquad
            t\sim N\left(0,\frac{4a\rho}{\sqrt\eta}Q_z^{-1}\right).
                                                               \tag{19}
\]
The radial density has the more precise equivalent
\[
 q_{\beta_N,z}(N\eta)\sim
 C_0 a^{r/2}\eta^{-(r+2)/4}N^{-(r+1)/2}(\log N)^{r/2}
                 e^{-\kappa_0\sqrt{N\eta}/\beta_N},      \tag{20}
\]
where
\[
 C_0=\frac{(2\pi)^{-(r+2)/2}}{2\mathcal M_\infty}
                  \sqrt{2\pi}\frac{(8\pi\rho)^{r/2}}{\sqrt{\det Q_z}}.
\]
This theorem is a critical one-neuron conditional limit; it does not
state a full-posterior fluctuation limit in the joint regime.
\end{theorem}
\begin{proof}
Substitute $u=\sqrt{\log N/N}\,t$ in the exact exponent (8).
On bounded marks, $\beta_N u\cdot\bar x=O((\log N)^{-1/2})$, and
Taylor expansion of $\log\cosh$ gives
\[
 D_{\beta_N}-\frac{\kappa_0\sqrt{N\eta}}{\beta_N}
        \longrightarrow \frac{v^2}{2}
                      +\frac{\sqrt\eta}{8a\rho}t^\top Q_z t.
\]
The imbalance correction to the subtracted constant is
$O(|v|/\beta_N+v^2/(\beta_N\sqrt N))$ and vanishes.
The Jacobian is asymptotic to
$\tfrac12 N^{-(r+1)/2}\eta^{-1/2}(\log N)^{r/2}$.
After this rescaling, (9) bounds the transverse exponent below by a
constant times $\min(\|t\|^2,\sqrt{\log N}\|t\|)$; the completed-square
bound controls $v$. These give integrable domination, including the
expanding cone boundary. Dominated convergence proves (19) and (20).
The exact exponential in (20) cannot in general be replaced by
$N^{-a\kappa_0\sqrt\eta}$ in an asymptotic equivalent: the assumption
$a_N\to a$ does not imply $(a_N-a)\log N\to0$.
\end{proof}

The one-neuron crossover companion identifies $\beta\asymp\sqrt c$
as the scale at which the transverse mark law changes from its smooth
Gaussian form to the ReLU corner law. At $c\asymp N$, Theorem~\ref{thm:main}
shows that the polynomial rare-count exponent already changes at
$\beta\asymp\sqrt N/\log N$. The logarithmic separation concerns two
different observables. It does not identify an additional thermodynamic
phase or assert uniformity in increasing input dimension.

\section{The first Gaussian-mark correction is logarithmic}
The one-neuron limit can be sharpened when the temperature is specified
exactly. This also distinguishes a local Gaussian approximation error
from the polynomial posterior count probabilities.

\begin{corollary}[A nonzero logarithmic mark error]\label{cor:mark-rate}
Fix $a,\eta>0$ and set $\beta_N=\sqrt N/(a\log N)$ exactly. Let $G$
be the Gaussian law in (19), and let $P_N^{\rm mark}$ be the conditional
mark law there. Define
\[
 Q_4(t)=\frac{\sqrt\eta}{192a^3\rho M}
                         \sum_{x:z\cdot x=0}(t\cdot\bar x)^4,
 \qquad \mu_4=\E_G Q_4=\frac{r^2\rho}{4a\sqrt\eta\,p_0}.
                                                               \tag{21}
\]
Then, extending its density by zero outside the cone,
\[
 \left\|\log N\left(\frac{dP_N^{\rm mark}}{dG}-1\right)
                           -(Q_4-\mu_4)\right\|_{L^1(G)}\to0.
                                                               \tag{22}
\]
In particular
\[
 \|P_N^{\rm mark}-G\|_{\rm TV}
       \sim\frac{\E_G|Q_4-\mu_4|}{2\log N}>0.
                                                               \tag{23}
\]
The ratio of the radial density to the right side of (20) equals
$1+\mu_4/\log N+o(1/\log N)$.
\end{corollary}
\begin{proof}
The next term of $\log\cosh u=u^2/2-u^4/12+O(u^6)$ gives the
unnormalized mark density relative to $G$ as
$1+Q_4(t)/\log N+o_{L^1(G)}(1/\log N)$.
For a detailed domination argument, first restrict
$|v|+\|t\|\le(\log N)^\epsilon$ with a sufficiently small fixed
$\epsilon>0$. The sixth-order remainder is bounded by
$C(\log N)^{-2}\|t\|^6$, and exponentiating the fourth-order term
costs $C(\log N)^{-2}(1+\|t\|^8)$ under a slightly weaker Gaussian
density. Balance, Jacobian, and nonzero-layer errors are smaller than
every fixed inverse power of $\log N$. On the complement, the
min-quadratic/linear bound used in Theorem~\ref{thm:marks} gives
$Ce^{-c(\log N)^{2\epsilon}}$ integrated mass for both the exact
unnormalized marks and the Gaussian polynomial terms. This proves
the claimed $L^1$ expansion. Normalization subtracts $\mu_4$.

For every zero-layer $x$,
$\bar x^\top Q_z^{-1}\bar x=r/p_0$ on the transverse space.
The fourth Gaussian moment is therefore
$3[4a\rho r/(\sqrt\eta p_0)]^2$. There are $Mp_0$ such rows, giving
(21). The nonconstant quartic polynomial has positive variance under
the nondegenerate Gaussian, proving strict positivity in (23).
Finally $\mathcal M_\beta-\mathcal M_\infty=o(1/\beta)$, as shown
in the next proof. At the present temperature this is
$o(1/\log N)$, so replacing the normalizer by $\mathcal M_\infty$
in the radial prefactor does not affect its displayed correction.
\end{proof}

\section{A singular covariance correction in six input dimensions}
Pointwise softplus approximation also hides a distinct sensitivity of
the critical second moments. The following statement concerns one
critical neuron, with inverse temperature tending to infinity; it is
separate from the joint posterior limit in Theorem~\ref{thm:main}.

\begin{theorem}[The critical covariance remembers the large-neuron tail]
\label{thm:covariance}
Fix $r=6$, and write
$\Sigma_\beta=\operatorname{Cov}_{Q_\beta}(\Phi_\beta)$.
On the transverse space let
\[
 I_\zeta=\int_{\R^6}e^{-V_\zeta(t)}\,dt,\qquad
 I_\infty=\int_{\R^6}e^{-V_\infty(t)}\,dt,
 \quad V_\infty(t)=\frac1{2\rho M}\sum_{z\cdot x=0}|t\cdot\bar x|.
\]
Then
\[
 \beta(\Sigma_\infty-\Sigma_\beta)
        \longrightarrow C_*\sum_{z\in X} f_zf_z^\top,
 \qquad
 C_* =\frac{(2\pi)^{-7/2}}{\mathcal M_\infty}
                         \int_0^\infty(I_\infty-I_\zeta)\,d\zeta
                         \in(0,\infty).                \tag{24}
\]
The limiting matrix has rank $38$. Its Walsh eigenvalues, by degree
$0,1,\ldots,6$, are $C_*$ times
\[
                     (600,\ 512/3,\ 24,\ 0,\ 8/3,\ 0,\ 24).
                                                               \tag{25}
\]
Thus the leading covariance change is positive on the critical-ray
span and zero on its orthogonal complement. No finite-$\beta$ matrix
ordering is asserted.
\end{theorem}
\begin{proof}
For fixed $\zeta>0$, write
\[
 \Delta V_\zeta(t)=V_\zeta(t)-V_\infty(t)
   =\frac1{\zeta\rho M}\sum_{z\cdot x=0}
                        \log(1+e^{-\zeta|t\cdot\bar x|})\ge0.
                                                               \tag{26}
\]
For any $q$ in a fixed compact positive interval, coarea on each
nonzero linear functional $t\mapsto t\cdot\bar x$ gives
\[
 \int e^{-qV_\infty(t)}(1-e^{-q\Delta V_\zeta(t)})\,dt
                              \le C\min\{1,\zeta^{-2}\}.
                                                               \tag{27}
\]
For $\zeta\ge1$ use $1-e^{-u}\le u$: the explicit $\zeta^{-1}$
in (26) and the one-dimensional integral of
$\log(1+e^{-\zeta|s|})$ each contribute a factor $\zeta^{-1}$.
The transverse hyperplane integral is uniformly bounded because
$V_\infty$ is a norm. For $\zeta<1$ the untilted integral bounds it.
The same estimate holds with any fixed polynomial in $\|t\|$ inserted.
In particular $0<I_\infty-I_\zeta\le C\min(1,\zeta^{-2})$;
the constant in (24) is finite and positive.

We first remove lower-order differences. Almost everywhere in parameter
space,
$\beta(\phi_\beta(U\cdot\bar x)-\phi_\infty(U\cdot\bar x))\to0$,
and the absolute value is at most $\log2$. Also
\[
 |e^u-e^v|\le |u-v|(e^u+e^v).
\]
The common cone/ReLU and off-cone Gaussian majorants in
Lemma~\ref{lem:corner}, with finite first and third parameter moments,
therefore imply
\[
 \beta(\mathcal M_\beta-\mathcal M_\infty)\to0,
 \qquad
 \beta(m_\beta-m_\infty)\to0.                         \tag{28}
\]
For the second-moment integrand, changing $\Phi_\beta\Phi_\beta^\top$
to $\Phi_\infty\Phi_\infty^\top$ under the soft critical density
also gives $o(\beta^{-1})$: its $\beta$-scaled difference is bounded
by $C(1+\|\theta\|^3)$ and tends to zero almost everywhere.
Off fixed critical cones, all remaining density differences have a
Gaussian majorant after multiplication by the required polynomials,
so their $\beta$-scaled contribution tends to zero as well.

On one large-$c$ critical cone, the exact ReLU output is
$cf_z+e$ with $\|e\|\le C\|t\|$. Replacing its tensor square by
$c^2f_zf_z^\top$ leaves $O(c\|t\|+\|t\|^2)$.
The mean-value density bound just used supplies one factor
$|a|\le C\sqrt c$ after multiplication by $\beta$. Against the
critical density $c^{-7/2}e^{-c_1(v^2+\|t\|)}$, this remainder is
integrable: its largest radial power is $c^{-2}$. Dominated convergence
makes this replacement another $o(\beta^{-1})$ error.

The exact exponent difference on the cone is
$q\Delta V_{\beta/\sqrt c}(t)$ plus an error bounded by
$C\sqrt c\,\beta^{-1}e^{-c_2\beta\sqrt c}$. The latter contributes
$o(\beta^{-1})$ even with the factor $c^2$. We are left with the
radial difference controlled, by (27), by
\[
                         Cc^{2-7/2}\min(1,c/\beta^2).   \tag{29}
\]
Put $c=\beta^2 y$. After multiplication by $\beta$, its integrable
majorant is $Cy^{-3/2}\min(1,y)$ on $(0,\infty)$.
At each fixed $y>0$, $q\to1$, the balance Jacobian tends to
$\tfrac12c^{-7/2}$, the physical transverse Gaussian factor tends
to one, and the expanding cone restrictions disappear. Integrating
$v$ gives $\sqrt{2\pi}$. Dominated convergence yields, per ray,
\[
 \frac{(2\pi)^{-7/2}}2 f_zf_z^\top
            \int_0^\infty y^{-3/2}(I_\infty-I_{1/\sqrt y})\,dy.
\]
The substitution $\zeta=1/\sqrt y$ changes the last integral to
$2\int_0^\infty(I_\infty-I_\zeta)\,d\zeta$.
Divide by $\mathcal M_\infty$, sum the finitely many rays, and use
(28) to normalize and center the second moments. This proves (24).

Finally the finite binomial sum gives the Walsh coefficients of the
six-bit ReLU sum, by degree, as
\[
                 (15/16,1/2,3/16,0,-1/16,0,3/16).
\]
Coordinate-flip symmetry diagonalizes $\sum_zf_zf_z^\top$; its
degree-$j$ eigenvalue is $M^2/6$ times the square of that coefficient.
This gives (25). The positive eigenspaces have total dimension
$1+6+15+15+1=38$.
\end{proof}

For numerical evaluation, a second positive-integral formula provides
an independent check. Homogeneity and Tonelli's theorem give
\[
 \int_0^\infty(I_\infty-I_\zeta)\,d\zeta
 =\Gamma(5)\int_{\R^6}
       [V_\infty(t)^{-5}-V_1(t)^{-5}]\,dt.             \tag{30}
\]
Indeed set $s=\zeta t$ inside $I_\zeta$ and then integrate
$\zeta^{-6}(e^{-V_\infty(s)/\zeta}-e^{-V_1(s)/\zeta})$.
Both integrands are nonnegative. Numerical values from these formulas
and paired critical-law integrations are Monte Carlo diagnostics, not
certified enclosures. The dimension-six conclusion is not asserted in
higher dimensions: the integrability balance in (29) changes with $r$.

\begin{corollary}[The highest finite integer output moment is singular]
Fix even $r\ge6$ and put $k=(r-2)/2$. Define $V_\zeta$ on the
$r$-dimensional transverse space as in (7), and let
$I_{\zeta,r}=\int_{\R^r}e^{-V_\zeta(t)}dt$.
Then the raw moment tensors obey
\[
 \beta\left(\E_{Q_\infty}\Phi_\infty^{\otimes k}
             -\E_{Q_\beta}\Phi_\beta^{\otimes k}\right)
 \longrightarrow C_r\sum_{z\in X}f_z^{\otimes k},\qquad
 C_r=\frac{(2\pi)^{-(r+1)/2}}{\mathcal M_\infty}
                 \int_0^\infty(I_{\infty,r}-I_{\zeta,r})d\zeta>0.
                                                               \tag{30a}
\]
The same limit holds for centered $k$th moment tensors. The limiting
tensor is nonzero: contracting each index against $\chi/M$ gives
$MC_r\rho^k>0$. The order $k$ is the largest integer for which the
absolute output moment is finite under $Q_\infty$.
\end{corollary}
\begin{proof}
The radial exponent is $\nu=(r+1)/2=k+3/2$, so the leading moment
density is $c^{k-\nu}=c^{-3/2}$. The argument for (27) works in every
fixed transverse dimension. The exact same substitution $c=\beta^2y$
therefore gives the integrable dominating function in (29) and the
constant in (30a). Replacing $\Phi_\infty^{\otimes k}$ by
$c^kf_z^{\otimes k}$ leaves at most a polynomial in the marks times
$c^{k-1}$; after the mean-value density bound, the largest radial power
is again $c^{k-1+1/2-\nu}=c^{-2}$. Changing the output itself has a
$\beta$-scaled bound of parameter degree $2k-1<r-1$. Both errors are
$o(1)$ after multiplication by $\beta$, by the same dominated-convergence
argument. Lower raw moment differences and mean differences are also
$o(\beta^{-1})$, because the corresponding density derivative needs
at most $2k-1$ parameter powers. Expanding a centered tensor proves the
centered statement. Finally the cone density is comparable to
$c^{-\nu}$ on bounded marks and the output norm is comparable to $c$;
its absolute moment is finite exactly for orders below
$\nu-1=k+1/2$. This proves the integer-order assertion.
\end{proof}

\section{The critical normalizer has a regular second-order expansion}
The normalizer in (2)--(3) contributes a larger correction than the
polynomial count cost at the logarithmic smoothing scale. It can be
expanded explicitly, provided its second-order term is retained.
Write $\varphi_j$ for the standard Gaussian density in $j$ dimensions,
and $G(U)=M^{-1}\sum_x\chi(x)\max(U\cdot\bar x,0)$.

\begin{theorem}[A boundary formula for the normalizer]\label{thm:normalizer}
For every fixed even $r\ge6$, there is a finite constant $A_r<0$
such that
\[
          \mathcal M_\beta=\mathcal M_\infty
                                  +\frac{A_r}{\beta^2}+O(\beta^{-3}).
                                                               \tag{31}
\]
One formula, with $b=-\sum_iw_i$, is
\[
 A_r=\frac{\pi^2}{6\rho^2}\int_{\R^r}
       \varphi_r(w)\varphi_1(b)G(w,b)
                 e^{G(w,b)^2/(2\rho^2)}\,dw.             \tag{32}
\]
The ordering of $(w,b)$ in this formula is only notational: the
preactivation is $b+w\cdot x$. The strict sign is proved after the
boundary expansion using a positive Fourier representation.
\end{theorem}
\begin{proof}
The softplus increment is
$d_\beta(s)=\phi_\beta(s)-\phi_\infty(s)
=\beta^{-1}\log(1+e^{-\beta|s|})$.
Its two elementary integrals are
\[
 \int_\R d_\beta(s)\,ds=\frac{\pi^2}{6\beta^2},\qquad
 \int_\R |s|d_\beta(s)\,ds=O(\beta^{-3}),\qquad
 \int_\R d_\beta(s)^2\,ds=O(\beta^{-3}).               \tag{33}
\]
The first follows by integrating the alternating exponential series for
$\log(1+e^{-u})$; the other two follow by scaling and exponential decay.

We record a uniform slice bound to justify the expansions at the
critical tilt. The sharp biased ReLU coefficient inequality gives
\[
             |G(U)|/\rho\le\sqrt r\min_i|w_i|.
\]
The same bound holds for $G_\beta$: softplus is ReLU averaged over a
centered logistic bias perturbation, and the inequality is uniform in
bias. It also holds for every convex combination of $G$ and $G_\beta$.
On a slice $b+w\cdot x=s$, integrate in $b$ and bound its Gaussian
density, or its first derivative, by a constant. Partition the incoming
weights by their signs and their minimum absolute coordinate $t$;
write the other magnitudes as $t+e_i$, $e_i\ge0$. Completing the
Gaussian readout square bounds its $p$th absolute moment by
$C_p(1+t)^p$. Integration over the $r-1$ excesses is bounded by
$C(1+t)^{-(r-1)}$. Consequently the slice integrals with readout moments
$p=1,2$ are uniformly finite in $s$, because
\[
                  \int_0^\infty(1+t)^{p-(r-1)}\,dt<\infty.
                                                               \tag{34}
\]
The same argument applies to the absolute derivative of the bias
Gaussian density.

Put $\delta G=M^{-1}\sum_x\chi(x)d_\beta(U\cdot\bar x)$.
Taylor expansion of $e^{a(G+\delta G)/\rho}$ about $G$ has integrated
remainder bounded by
\[
 C\sum_x\int a^2d_\beta(U\cdot\bar x)^2
           [e^{aG/\rho}+e^{aG_\beta/\rho}]\,P_0(d\theta)
                                                    =O(\beta^{-3}),
\]
using (33)--(34) and the convex-combination bound just established.
For each $x$ the linear term is
$\int d_\beta(s)F_x(s)\,ds$, where $F_x(s)$ is the Gaussian slice
integral of $a\chi(x)e^{aG/\rho}/(M\rho)$ at $b=s-w\cdot x$.
This function is uniformly Lipschitz. Differentiating almost everywhere
uses the first derivative of the bias Gaussian and
$|\partial_bG|\le1$; the resulting readout powers are at most two,
so (34) bounds the derivative uniformly. Therefore (33) gives
\[
       \int d_\beta(s)F_x(s)\,ds
                =\frac{\pi^2}{6\beta^2}F_x(0)+O(\beta^{-3}).
\]
Sum over the finite cube. Integrating the Gaussian readout in $F_x(0)$
replaces $a e^{aG/\rho}$ by $(G/\rho)e^{G^2/(2\rho^2)}$.
Finally coordinate flips send each $x$ slice to the all-positive input
slice and cancel its factor $\chi(x)$. The $M$ equal contributions
cancel the factor $1/M$, giving (32).
\end{proof}

\section{Fourier moments and normalized initialization spectra}
\begin{theorem}[Smoothing orders all even parity-correlation moments]
\label{thm:moment-order}
Let $b,w_1,\ldots,w_r$ be independent standard Gaussians and put
$G_\beta(b,w)=M^{-1}\sum_x\chi(x)\phi_\beta(b+w\cdot x)$.
For every even $r\ge2$ and integer $j\ge1$,
$\E G_\beta^{2j}$ is strictly increasing in
$\beta\in(0,\infty]$. More generally the same weak ordering holds
for any two centered ReLU bias smoothings with finite first moments whose characteristic
multipliers satisfy $0\le\psi_1(t)\le\psi_2(t)$ for all $t>0$.
Consequently, for even $r\ge6$,
\[
       1<\mathcal M_{\beta_1}<\mathcal M_{\beta_2}
                         <\mathcal M_\infty
            \qquad(0<\beta_1<\beta_2<\infty).          \tag{M1}
\]
Both endpoint limits are exact: $\mathcal M_\beta\to1$ as
$\beta\downarrow0$, and $\mathcal M_\beta\to\mathcal M_\infty$
as $\beta\to\infty$.
\end{theorem}
\begin{proof}
Softplus is ReLU averaged over a centered logistic bias with scale
$1/\beta$. Its characteristic function is
\[
        \psi_\beta(t)=\frac{\pi t/\beta}{\sinh(\pi t/\beta)},
        \qquad \psi_\infty(t)=1.
\]
For completeness, substituting $u=e^{-z}$ in the logistic-density
Fourier integral gives
$\int_0^\infty u^{-it}(1+u)^{-2}du
=\Gamma(1-it)\Gamma(1+it)=\pi t/\sinh(\pi t)$.
The ReLU cosine integral, after the affine terms cancel against parity,
therefore gives
\[
 G_\beta(b,w)=\frac{(-1)^{r/2+1}}\pi
       \int_0^\infty \psi_\beta(t)\cos(bt)
                    \prod_{i=1}^r\sin(w_it)\,\frac{dt}{t^2}.
                                                               \tag{M2}
\]
For $\mathbf t=(t_1,\ldots,t_{2j})$, let
\[
 C_j(\mathbf t)=\E_Z\prod_{h=1}^{2j}\cos(Zt_h),\qquad
 S_j(\mathbf t)=\E_Z\prod_{h=1}^{2j}\sin(Zt_h),
                      \quad Z\sim N(0,1).
\]
Expanding the cosines shows
$C_j=2^{-2j}\sum_{\varepsilon\in\{-1,1\}^{2j}}
\exp[-(\sum_h\varepsilon_ht_h)^2/2]>0$.
Independence of the incoming Gaussian coordinates gives
\[
 \E G_\beta^{2j}=\pi^{-2j}\int_{(0,\infty)^{2j}}
       C_j(\mathbf t)S_j(\mathbf t)^r
                    \prod_{h=1}^{2j}\frac{\psi_\beta(t_h)}{t_h^2}
                         \,d\mathbf t.                 \tag{M3}
\]
Fubini is justified before averaging: for fixed weights the absolute
integral in (M2) is bounded by a polynomial in $1+\|w\|$, using
$|\sin(w_it)|\le\min(1,|w_i|t)$ near zero and $t^{-2}$ at infinity.
Its $2j$th power has finite Gaussian expectation.
Since $r$ is even, the entire integrand in (M3) is nonnegative.
The multiplier increases strictly with $\beta$ at every $t>0$.
Moreover $S_j$ is nonzero on an open set near positive equal small
$t_h$, since its leading term is
$\E Z^{2j}\prod_h t_h>0$. This proves strict moment ordering.

Integrating the independent Gaussian readout gives
$\mathcal M_\beta=\E\exp(G_\beta^2/(2\rho^2))$.
Expand this exponential into its nonnegative moment series and use
(M3). Finiteness of $\mathcal M_\infty$ gives domination for the
series and the integrals. Pointwise multiplier convergence to zero
or one proves both endpoint limits; the $j=1$ term proves strictness.
\end{proof}

\begin{corollary}[A centered random-feature kernel ordering]
Let $d\ge1$ and initialize the bias and $d$ incoming weights as
independent standard Gaussians. On $\{-1,1\}^d$ center each feature
over the input cube:
\[
 h_\beta(U,x)=\phi_\beta(U\cdot\bar x)
                         -\E_z\phi_\beta(U\cdot\bar z),\qquad
 K_\beta^c(x,y)=\E_U h_\beta(U,x)h_\beta(U,y).
\]
If $0<\beta_1<\beta_2\le\infty$, then
$K_{\beta_2}^c-K_{\beta_1}^c$ is positive definite on the mean-zero
subspace and is zero on constants. Every nonconstant Walsh eigenvalue
strictly increases as smoothing is removed.
Consequently every nonzero mean-zero target has a strictly smaller
RKHS norm under $K_{\beta_2}^c$ than under $K_{\beta_1}^c$.
This is an initialization-kernel statement, with no assertion of
trained-network or sample-complexity monotonicity.
\end{corollary}
\begin{proof}
Coordinate-flip invariance makes the kernel diagonal in the Walsh
basis. With the normalized counting measure, its eigenvalue for
$S\ne\varnothing$ is $\E\widehat\phi_\beta(S)^2$.
For degree $p=|S|\ge2$, the integral (M2) has sine factors on $S$
and cosine factors off $S$. The bias factor is cosine when $p$ is
even and sine when $p$ is odd. Thus its $2j$th moment has nonnegative
integrand proportional to
\[
 \begin{cases}
 S_j^p C_j^{d-p+1}\prod_h\psi_\beta(t_h)/t_h^2,&p\text{ even},\\
 S_j^{p+1}C_j^{d-p}\prod_h\psi_\beta(t_h)/t_h^2,&p\text{ odd}.
 \end{cases}
\]
Both sine exponents are even; the isotropic Gaussian bias is used
in the odd case. This proves strict moment ordering for these
coefficients, in particular their second moments.
For degree one there is also the affine coefficient $w_i/2$.
The nonlinear coefficient is odd in the bias and has zero covariance
with $w_i/2$, so its eigenvalue is $1/4$ plus the same positive
integral with factors $S_1^2C_1^{d-1}$. It too strictly increases.
All nonconstant eigenvalues are positive. The finite-cube RKHS norm
is $\sum_{S\ne\varnothing}|\widehat f(S)|^2/\lambda_S$,
which proves the final assertion.
\end{proof}


Activation-dependent spectral regularization is studied, in a different
Radon-transform kernel setting, by Sahs et al.~\cite{sahs}.
Mehrpanah et al.~\cite{mehrpanah} study Gaussian activation smoothing
and NTK spectral tails. The next result is an exact finite-cube,
input-centered random-feature statement with Gaussian bias: it orders
all degree ratios and all normalized degree tails for logistic ReLU
smoothing. No general novelty claim about smoothing and spectral bias
is intended.

\begin{theorem}[Sharpening shifts normalized variance toward higher degrees]
\label{thm:degree-order}
In the preceding centered-kernel model, let $d\ge2$, and let
$\lambda_p(\beta)$ be the eigenvalue of any degree-$p$ Walsh character.
For $1\le p<q\le d$, the ratio
$\lambda_q(\beta)/\lambda_p(\beta)$ is nondecreasing in
$\beta\in(0,\infty]$. It is strictly increasing except for the
identically equal pairs $p=2j-1$, $q=2j$, $j\ge2$.
Define the average centered feature variance and its degree distribution by
\[
 T_\beta=\sum_{p=1}^d\binom dp\lambda_p(\beta),\qquad
 \pi_p(\beta)=\frac{\binom dp\lambda_p(\beta)}{T_\beta}.
                                                               \tag{K1}
\]
Every upper tail $\sum_{p=k}^d\pi_p(\beta)$, $2\le k\le d$,
strictly increases with $\beta$. In particular, under the
unit-average-variance kernel $K_\beta^c/T_\beta$, the squared RKHS
norm of full parity, $T_\beta/\lambda_d(\beta)$, strictly decreases.
These conclusions concern the initialization kernel only.
\end{theorem}
\begin{proof}
Write, for $t,s>0$,
\[
 C=e^{-(t^2+s^2)/2}\cosh(ts),\qquad
 S=e^{-(t^2+s^2)/2}\sinh(ts).
\]
\[
 L_j(\beta)=\frac1{\pi^2}\int_0^\infty\!\int_0^\infty
 \frac{\psi_\beta(t)\psi_\beta(s)}{t^2s^2}
       S^{2j}C^{d+1-2j}\,dt\,ds.
                                                               \tag{K2}
\]
Here $1\le j\le\lfloor(d+1)/2\rfloor$. The preceding proof gives
$\lambda_1=1/4+L_1$ and
$\lambda_p=L_{\lceil p/2\rceil}$ for $p\ge2$.
All these integrals are finite and strictly positive, including at
$\beta=\infty$.

Fix $i<j$ and $0<\beta_1<\beta_2\le\infty$. Normalize the integrand
of $L_i(\beta_1)$ to a probability density on $(0,\infty)^2$.
Its log density has mixed derivative, with $u=ts$,
\[
 2i\{\coth u-u\operatorname{csch}^2u\}
 +(d+1-2i)\{\tanh u+u\operatorname{sech}^2u\}>0.       \tag{K3}
\]
The first brace is positive since $\sinh u\cosh u>u$.
Consequently the conditional density of $s$ given $t$ is increasing
in $t$ in likelihood-ratio order. We spell out the elementary
association implication. For two coordinatewise increasing functions
$f,g$, conditional one-dimensional Chebyshev gives
$\operatorname{Cov}(f,g\mid t)\ge0$.
The conditional expectations of $f$ and $g$ increase with $t$,
because both their explicit dependence on $t$ and the conditional
law of $s$ increase. Another one-dimensional Chebyshev inequality
therefore gives $\operatorname{Cov}(f,g)\ge0$.
Here Chebyshev follows from
$2\operatorname{Cov}(v(X),w(X))
=\E[(v(X)-v(X'))(w(X)-w(X'))]\ge0$ for an independent copy $X'$.
Truncation extends the argument to the nonnegative integrable
functions used below.

Take
\[
 f(t,s)=\tanh(ts)^{2(j-i)},\qquad
 g(t,s)=\frac{\psi_{\beta_2}(t)\psi_{\beta_2}(s)}
                 {\psi_{\beta_1}(t)\psi_{\beta_1}(s)}.
\]
Both are strictly increasing in each coordinate. Indeed, with
$a_h=\pi/\beta_h$ and $a_2=0$ at infinity,
\[
 \frac d{dt}\log\frac{\psi_{\beta_2}(t)}{\psi_{\beta_1}(t)}
       =a_1\coth(a_1t)-a_2\coth(a_2t)>0,
\]
where the second term is interpreted as $1/t$ when $a_2=0$;
$u\coth u$ has positive derivative. Their conditional covariance
is strictly positive for every $t>0$, since the conditional density
is positive on the entire positive axis. Thus $\operatorname{Cov}(f,g)>0$.
The function $f$ is bounded, and $g$ is integrable because
$L_i(\beta_2)<\infty$. Reweighting by $g$ now proves
\[
             \frac{L_j(\beta_2)}{L_i(\beta_2)}
                >\frac{L_j(\beta_1)}{L_i(\beta_1)}.     \tag{K4}
\]
For comparisons with degree one, use the additional strictly
increasing factor $L_1/(1/4+L_1)$, whose monotonicity follows from
the preceding corollary. This proves all eigenvalue-ratio claims.

Finally the ratios $\pi_p(\beta_2)/\pi_p(\beta_1)$ are
nondecreasing in $p$. For each cut $2\le k\le d$, every cross term
with $p<k\le q$ satisfies
$\lambda_q(\beta_2)\lambda_p(\beta_1)
\ge\lambda_q(\beta_1)\lambda_p(\beta_2)$, and the term
$p=1$, $q=d$ is strict. Multiplying by the positive binomial
multiplicities and summing proves the strict upper-tail claim.
The case $k=d$ proves the normalized parity-norm statement.
\end{proof}
The same argument applies to two centered symmetric ReLU bias smoothings
with finite first moments and strictly positive characteristic multipliers $\psi_1,\psi_2$,
provided $\psi_2/\psi_1$ is strictly increasing on $(0,\infty)$ and
the base two-variable density in (K2) is integrable. Its mixed log
derivative is still (K3), since the multipliers are separate functions
of $t$ and $s$. The affine degree-one comparison additionally uses
$\psi_2\ge\psi_1$. In particular all conclusions hold when reducing
the variance of independent Gaussian bias smoothing, whose multiplier
is $e^{-vt^2/2}$.


\begin{theorem}[The same degree ordering for the initialization tangent kernel]
\label{thm:tangent-order}
Keep the isotropic Gaussian incoming initialization above, and give the
readout independent variance one. The centered two-layer initialization
neural tangent kernel, when both layers are trainable, is
\[
 K_\beta^{\rm tan,c}(x,y)
 =\E_U\{h_\beta(U,x)h_\beta(U,y)
            +\nabla_Uh_\beta(U,x)\cdot\nabla_Uh_\beta(U,y)\}.
                                                               \tag{K5}
\]
All conclusions of Theorem~\ref{thm:degree-order} hold for this kernel:
higher-to-lower Walsh eigenvalue ratios increase with $\beta$, with
exactly the stated paired-degree exceptions; every normalized upper
degree tail strictly increases; and normalized full-parity RKHS cost
strictly decreases. The same is true for Gaussian ReLU bias smoothing.
The same ordering holds after multiplying the feature term and gradient
term by any fixed nonnegative constants, not both zero.
This concerns the initialization tangent kernel and its fixed-kernel
linearization, not an evolving finite-network kernel.
\end{theorem}
\begin{proof}
Put $q=d+1$, $u=ts$, and for $1\le j\le\lfloor q/2\rfloor$ define
\[
 H_j(u)=\sinh(u)^{2j}\cosh(u)^{q-2j},\qquad
 F_j(u)=H_j(u)+uH_j'(u).
\]
Differentiating the Fourier coefficient under its integral shows that
the tangent eigenvalue is
\[
 \nu_p(\beta)=\frac1{\pi^2}\int_0^\infty\!\int_0^\infty
 \frac{\psi_\beta(t)\psi_\beta(s)}{t^2s^2}
 e^{-q(t^2+s^2)/2}F_{\lceil p/2\rceil}(ts)\,dt\,ds
                +\tfrac12\one_{\{p=1\}}.              \tag{K6}
\]
Indeed differentiation of each sine factor contributes $ts\cosh(ts)$
in place of $\sinh(ts)$ in the squared-coefficient integral, and
differentiation of each cosine factor contributes $ts\sinh(ts)$
in place of $\cosh(ts)$. Summing gives $uH_j'$.
The affine coefficient contributes $1/4$ to each of the feature and
incoming-gradient variances at degree one; the cross terms vanish
by oddness in the bias. At finite $\beta$, exponential multiplier decay
justifies differentiation under the Fourier integral and Fubini; near
zero the sine factors give an integrable bound with finite Gaussian
weight moments. As $\beta\to\infty$, the activation derivatives
converge almost surely and are bounded by one, while the centered
features converge in Gaussian $L^2$. Thus the tangent variances converge
by dominated convergence. The positive integrals increase to the
ReLU expression by monotone convergence, proving (K6) at that endpoint
as well. The Gaussian-smoothing case has the same justification.

We need two positivity facts about $F_j$. First, it has a nonnegative
power series. For any such series $F(u)=\sum_n a_nu^n>0$,
\[
 \frac{d}{du}\{u(\log F)'(u)\}
     =\frac1u\operatorname{Var}_{\omega_u}(n)\ge0,
 \qquad \omega_u(n)=a_nu^n/F(u).                       \tag{K7}
\]
Consequently the integrand in (K6), normalized at any fixed scale,
has nonnegative mixed log derivative. Its conditional distributions
are ordered in likelihood ratio, so the association argument above applies.

Second, $F_j(u)/F_i(u)$ strictly increases in $u>0$ for $i<j$.
Here is a finite-matrix proof. The adjacency matrix of the $q$-cube,
restricted to normalized Hamming-level indicators, is the symmetric
tridiagonal matrix $B$ with
$B_{p,p+1}=\sqrt{(p+1)(q-p)}$. On the even levels let $T$ be the
restriction of $B^2$. It is nonnegative, tridiagonal, and positive
semidefinite. Every two-by-two minor of $T$ is nonnegative:
the only potentially negative ones are adjacent principal minors,
and these are nonnegative by positive semidefiniteness.
Multiplication by such a matrix preserves likelihood-ratio order
of nonnegative row vectors. Explicitly, for $j<k$,
\[
 (vT)_j(wT)_k-(vT)_k(wT)_j
 =\sum_{a<b}(v_aw_b-v_bw_a)
                      (T_{aj}T_{bk}-T_{ak}T_{bj})\ge0
\]
whenever $v_aw_b\ge v_bw_a$ for $a<b$.
Since the row vector concentrated at level zero precedes its image
under $T$ in this order, induction orders the rows $e_0T^n$ as $n$
increases.

The cube exponential identity gives
\[
 (e^{uB})_{0,2j}=\sqrt{\binom q{2j}}\,H_j(u),\qquad
 [u^{2n}]H_j(u)=
       \frac{(T^n)_{0,2j}}{(2n)!\sqrt{\binom q{2j}}}.     \tag{K8}
\]
Thus the ratio of the degree-$2n$ coefficients of $H_j$ and $H_i$
is nondecreasing in $n$, interpreted by cross products when coefficients
vanish. Passing from $H$ to $F=H+uH'$ multiplies both coefficients
by the same positive factor $2n+1$, preserving the order.
Pairwise expansion of the derivative of $F_j/F_i$ now gives
nonnegative terms. It is strictly positive: the coefficient at $n=i$
is positive for $H_i$ and zero for $H_j$, whereas at $n=j$ both
are positive. This proves the second fact.

The ratio $F_j(ts)/F_i(ts)$ is therefore strictly increasing in both
coordinates. Reweight by the strictly increasing multiplier ratio
$\psi_{\beta_2}(t)\psi_{\beta_2}(s)/
 [\psi_{\beta_1}(t)\psi_{\beta_1}(s)]$ and apply conditional Chebyshev
as in (K4). Integrability follows from the finite tangent variances;
if the ratio is unbounded, truncate first and use monotone convergence.
This proves strict ratios for the integral parts of (K6).
Each integral also increases strictly with $\beta$. Handle the extra
$1/2$ at degree one exactly as the earlier $1/4$, then sum cross
products to obtain all normalized upper-tail assertions.
For nonnegative term weights $c_0,c_1$, replace $F_j$ by
$c_0H_j+c_1uH_j'$. Its coefficient multiplier is $c_0+2nc_1>0$
on every nonzero coefficient under consideration. All arguments persist,
and the affine variance becomes $(c_0+c_1)/4$.
\end{proof}
For example, full-cube squared-loss gradient flow with this
\emph{fixed}, unit-average-variance kernel, zero initial prediction,
and parity target has residual proportional to
$\exp[-t\nu_d/\sum_p\binom dp\nu_p]$ (with the usual half-squared-loss
normalization). The theorem strictly improves this rate as smoothing
is removed. It does not assert the same comparison for finite-network
gradient descent or partially observed training.

\section{Complete monotonicity and higher-order boundary responses}
\begin{corollary}[Complete monotonicity in the smoothing variance]
For fixed even $r\ge6$, set $\mathfrak M(u)=\mathcal M_{1/\sqrt u}$
for $u>0$ and $\mathfrak M(0)=\mathcal M_\infty$.
There is a finite positive measure $\mu_r$ on $[0,\infty)$ with
$\mu_r(\{0\})=1$ and total mass $\mathcal M_\infty$ such that
\[
             \mathfrak M(u)=\int_0^\infty e^{-u z}\,\mu_r(dz).
                                                               \tag{M6}
\]
In particular $(-1)^n\mathfrak M^{(n)}(u)>0$ for every integer
$n\ge1$ and $u>0$, and $\log\mathfrak M$ is convex.
\end{corollary}
\begin{proof}
The product formula for $\sinh$ gives
\[
 \psi_{1/\sqrt u}(t)=\prod_{n\ge1}(1+ut^2/n^2)^{-1}
       =\E\exp\left[-ut^2\sum_{n\ge1}E_n/n^2\right],
\]
where the $E_n$ are independent unit-rate exponentials.
The sum is finite almost surely because its expectation is
$\sum n^{-2}<\infty$, and is strictly positive.
Insert independent copies into (M3) and then the exponential moment
series for $\mathcal M$. All mixing weights are nonnegative and their
total mass at $u=0$ is finite. Tonelli gives (M6); only the constant
term of the series contributes an atom at zero. Differentiation at
$u>0$ is dominated by $z^ne^{-uz}$, bounded after a smaller exponential
factor, and gives strict alternating derivatives.
Cauchy--Schwarz under the tilted measure gives log-convexity.
\end{proof}

\begin{corollary}[The normalizer response is strictly negative]
The coefficient $A_r$ in (31)--(32) satisfies
\[
 -A_r\ge\frac{D_r}{2\rho^2}>0,\qquad
 D_r=\frac16\int_0^\infty\!\int_0^\infty
 \frac{t^2+s^2}{t^2s^2}
 e^{-(r+1)(t^2+s^2)/2}\cosh(ts)\sinh(ts)^r\,dt\,ds.
                                                               \tag{M4}
\]
\end{corollary}
\begin{proof}
For $j=1$,
$C_1=e^{-(t^2+s^2)/2}\cosh(ts)$ and
$S_1=e^{-(t^2+s^2)/2}\sinh(ts)$.
The expansion $\psi_\beta(t)=1-\pi^2t^2/(6\beta^2)+O(\beta^{-4}t^4)$
and dominated convergence in (M3) give
$\beta^2(\E G_\infty^2-\E G_\beta^2)\to D_r$.
One useful global domination is
$\beta^2(1-\psi_\beta(t)\psi_\beta(s))\le C(t^2+s^2)$.
The integral in (M4) is finite: near either axis the sine power
removes the denominator; at infinity the possible large contribution
lies near $t=s$, where it is bounded by
$C(t^{-2}+s^{-2})e^{-(r+1)(t-s)^2/2}$ away from the axes.
It is positive on the positive quadrant.
Every term of the exponential moment series has a nonnegative
ReLU-minus-softplus difference, so
$\mathcal M_\infty-\mathcal M_\beta
\ge(\E G_\infty^2-\E G_\beta^2)/(2\rho^2)$.
Use (31) and pass to the limit to prove (M4), completing the strict-sign
claim in Theorem~\ref{thm:normalizer}.
\end{proof}

The positive number in (M4) also has a finite radical formula,
\[
 D_r=-\frac{\pi}{6\,2^r}\sum_{j=0}^r(-1)^j\binom rj
             \sqrt{1-\left(\frac{r+1-2j}{r+1}\right)^2}.
                                                               \tag{M5}
\]
To verify it independently, replace logistic smoothing by centered
Gaussian smoothing of variance $v=\pi^2/(3\beta^2)$.
Its Fourier multiplier has the same second-order expansion and the
same domination used for (M4), so its limiting second-moment response
is $D_r$. The standard Gaussian ReLU covariance formula
\cite{cho-saul}, for equal marginal variances $q$ and covariance $c$, is
\[
 K(q,c)=\frac1{2\pi}
       [\sqrt{q^2-c^2}+c(\pi-\arccos(c/q))].
\]
After independent smoothing in the two factors, $q=r+1+v$ and
$c=r+1-2j$. Thus the parity second moment is
$2^{-r}\sum_j(-1)^j\binom rj K(r+1+v,r+1-2j)$.
Its right derivative at $v=0$ uses
$\partial_qK=\sqrt{1-(c/q)^2}/(2\pi)$, including the endpoint
derivative zero at $c=q$, and gives (M5).
For example $D_6/(2\rho^2)\approx0.2209004$, whereas the separate
boundary-integral estimate gives $-A_6\approx0.3969$.
These decimal evaluations are numerical, not interval enclosures.

\begin{theorem}[A positive third-order boundary correction]
For every fixed even $r\ge6$,
\begin{gather}
 \mathcal M_\beta=\mathcal M_\infty+
             A_r\beta^{-2}+J_r\beta^{-3}+o(\beta^{-3}),\notag\\
 J_r=\frac{\zeta(3)}{M\rho^2}\int_{\R^r}
  \varphi_r(w)\varphi_1(b)
  \left(1+\frac{G(w,b)^2}{\rho^2}\right)
                 e^{G(w,b)^2/(2\rho^2)}dw>0,\notag\\
 b=-\sum_iw_i.                                           \tag{M7}
\end{gather}
Thus the remainder order in (31) is sharp, with a positive leading
coefficient. Equivalently the completely monotone function in (M6)
has a finite first derivative but no finite second derivative at zero.
\end{theorem}
\begin{proof}
Retain one further order in the readout-space Taylor argument for (31).
The cubic remainder is $O(\beta^{-4})$: use
$|\delta G|^3\le C\sum_x d_\beta(U\cdot\bar x)^3$,
the single-slice bound (34) with readout power three, and
$\int d_\beta(s)^3ds=O(\beta^{-4})$.
Power three is integrable for $r\ge6$.

In the linear term write $c_x=\chi(x)/M$, and let $F_x(s)$ be
the slice function used above. Its derivative has one-sided limits
at zero. Almost every point of that slice lies on no other activation
hyperplane, so the only jump is the current ReLU kink:
\[
 F'_x(0+)-F'_x(0-)
   =\int_{b=-w\cdot x}\frac{a^2c_x^2}{\rho^2}
                         e^{aG/\rho}\,P_0(da,dw,db)/db.
\]
The notation on the right means evaluation of the Gaussian density
at that value of $b$, followed by integration in $a,w$.
Dominated convergence is justified by (34) with powers at most two.
Therefore
$F_x(s)=F_x(0)+\frac12[F'_x(0+)+F'_x(0-)]s
+\frac12[F'_x(0+)-F'_x(0-)]|s|+o(|s|)$.
Since $d_\beta$ is even, its odd contribution integrates to zero,
and
\[
 \int |s|d_\beta(s)ds=\frac{3\zeta(3)}{2\beta^3},\qquad
 \int d_\beta(s)^2ds=\frac{\zeta(3)}{2\beta^3}.          \tag{M8}
\]
The first integral follows from the alternating series for the logarithm;
the second from
$\int_0^1\log^2(1+x)\,dx/x=\zeta(3)/4$.
The linear term thus contributes $3\zeta(3)/4$ times each displayed
jump, and the quadratic self term contributes $\zeta(3)/4$ times it.

Quadratic terms from two distinct input hyperplanes are
$O(\beta^{-4})$. To justify this even at the critical tilt, impose
both slices, eliminate $b$, and use their nonzero difference to
eliminate one incoming coordinate. In minimum-coordinate variables,
if the eliminated coordinate is not the minimum, at most one excess
integral is lost. The remaining bound is
$C\int_0^\infty(1+t)^{2-(r-2)}dt<\infty$ for $r\ge6$.
If the difference equation fixes only the minimum coordinate,
its radial integral is replaced by a bounded evaluation of
$(1+t)^{2-(r-1)}$. Nonzero coefficients in the difference equation
are $\pm2$, so the coarea factors are uniformly bounded.
These estimates are uniform in the two slice values.
Multiply by the two integrals $\int d_\beta=O(\beta^{-2})$.

The two self contributions sum to $\zeta(3)$ times the jump.
Gaussian readout integration changes
$a^2e^{aG/\rho}$ to $(1+G^2/\rho^2)e^{G^2/(2\rho^2)}$.
Coordinate-flip symmetry makes all $M$ slices equal; their factors
$c_x^2=M^{-2}$ leave $1/M$. This proves (M7), including finiteness
and strict positivity. The positive $u^{3/2}$ term after the linear
term rules out a finite second derivative at $u=0$.
\end{proof}

\begin{corollary}[An explicit logarithm-squared evidence correction]
Under (1), the exact normalizer in Theorem~\ref{thm:main} can be replaced
by
\[
 e^{N\log\mathcal M_{\beta_N}}
 =\exp\left\{N\log\mathcal M_\infty+
                    \frac{A_rN}{\mathcal M_\infty\beta_N^2}+o(1)\right\}.
                                                               \tag{35}
\]
In particular, if $\beta_N=\sqrt N/(a\log N)$ exactly for fixed $a>0$,
the hard evidence is
\[
 \exp\left\{-N(s/\rho-\log\mathcal M_\infty)
           +\frac{A_ra^2}{\mathcal M_\infty}(\log N)^2\right\}
          N^{-C_H-a\kappa_0K\sqrt\alpha+o(1)},           \tag{36}
\]
and the soft analogue replaces $s/\rho$ by $g(s)$ and $C_H$ by $C_L$.
For merely convergent $a_N$, retain $N/\beta_N^2$ as in (35): substituting
$a^2(\log N)^2$ can leave an error larger than $\log N$.
\end{corollary}
\begin{proof}
Expand the logarithm in (31). Its remainder after multiplying by $N$
is $O(N\beta_N^{-3}+N\beta_N^{-4})=o(1)$ under (1).
Insert this in (2)--(3). The special exact temperature relation gives
(36).
\end{proof}

\section{Regular means and the higher-dimensional covariance response}
The singular covariance correction is dimension-sensitive. The first
critical moment remains regular even in dimension six, while the second
moment has a regular boundary response once the even dimension is at
least eight. These statements concern critical single-neuron laws.

\begin{theorem}[The critical mean has a second-order boundary response]
For every even $r\ge6$, put
\[
 B_r=\frac{\pi^2}{6}\int_{\R^r}\varphi_r(w)\varphi_1(b)
       \left(\frac{2G(w,b)}\rho+\frac{G(w,b)^3}{\rho^3}\right)
          e^{G(w,b)^2/(2\rho^2)}\,dw,\qquad b=-\sum_iw_i.
                                                               \tag{37}
\]
This integral is finite, and
\[
 m_\beta=m_\infty+
           \frac{B_r-m_\infty A_r}{\mathcal M_\infty\beta^2}
                         +o(\beta^{-2}).                \tag{38}
\]
No sign of the coefficient is asserted.
\end{theorem}
\begin{proof}
Integrating the Gaussian readout gives
$\mathcal M_\beta m_\beta=\E_U h(G_\beta(U))$, where
$h(g)=g^2e^{g^2/(2\rho^2)}/\rho$.
Use the fundamental theorem of calculus on the line from $G$ to
$G_\beta$, and substitute
$G_\beta-G=M^{-1}\sum_x\chi(x)d_\beta(U\cdot\bar x)$.
For each slice, the resulting integrand is the average of $h'(G_t)$,
$0\le t\le1$, times $\chi(x)/M$. The same critical minimum-coordinate
bound used in (34) applies to $G_t$. Since
\[
 h'(g)=\left(\frac{2g}\rho+\frac{g^3}{\rho^3}\right)
                                      e^{g^2/(2\rho^2)},
\]
the slice integrals are bounded by (34) with powers at most three;
$p=3$ is integrable for every $r\ge6$.
After $s=u/\beta$, the scaled increment measure is
$\beta^2d_\beta(s)\,ds=\log(1+e^{-|u|})\,du$.
The slice integrand tends pointwise to its unsmoothed value at $s=0$,
and the uniform slice bound gives dominated convergence in $u$ and
in the interpolation parameter. Thus the change in the unnormalized
mean is $B_r/\beta^2+o(\beta^{-2})$. Coordinate flips reduce all
slices to (37), since $h'$ is odd. Divide by (31) to obtain (38).
\end{proof}

Here is also an explicit, though less compact, formula for the
higher-dimensional covariance coefficient. Write
\[
 q_x(U)=\max\{U\cdot\bar x,0\},\qquad
 W(g)=(1+g^2/\rho^2)e^{g^2/(2\rho^2)},\qquad
 T_{xy}=\E_U[W(G(U))q_x(U)q_y(U)].
\]
For a function $H(U)$, let
\[
 \mathcal S_z(H)=\int_{\R^r}\varphi_r(w)\varphi_1(-w\cdot z)
                                      H(w,-w\cdot z)\,dw.
\]
This is the Gaussian slice integral with respect to the preactivation
coordinate $s=b+w\cdot z$; no surface-Jacobian factor is missing.

\begin{theorem}[Regular covariance response in even dimension at least eight]
Fix even $r\ge8$, and define the finite matrix
\[
 D_{xy}=\frac{\pi^2}{6}\sum_{z\in X}\mathcal S_z\left(
   \frac{\chi(z)}M W'(G)q_xq_y
       +W(G)[\one_{z=x}q_y+\one_{z=y}q_x]\right).        \tag{39}
\]
With $b_r=(B_r-m_\infty A_r)/\mathcal M_\infty$, one has
\[
 \beta^2(\Sigma_\beta-\Sigma_\infty)
 \longrightarrow \frac{D}{\mathcal M_\infty}
          -\frac{A_r T}{\mathcal M_\infty^2}
          -2m_\infty b_r\chi\chi^\top.                \tag{40}
\]
In particular the covariance change is $O(\beta^{-2})$ in these
fixed dimensions. No sign or nonzero-rank assertion is made for (40).
\end{theorem}
\begin{proof}
Gaussian integration in the readout shows that the unnormalized
second moment is
$T_{\beta,xy}=\E_U[W(G_\beta)\phi_\beta(U\cdot\bar x)
                                      \phi_\beta(U\cdot\bar y)]$.
Interpolate simultaneously between $G$ and $G_\beta$, and between
each ReLU value and its softplus value. Differentiation along this
finite line segment gives three terms: one with $W'(G_t)$ times
$\delta G$, and one for each changed activation. Substitute the
finite sum of increments for $\delta G$ and use the same scaled
approximate-identity argument as in (38).

For completeness, the domination now needs five incoming-weight
powers. The largest term is $W'(G_t)q_{t,x}q_{t,y}$, where
$W'(g)=(3g/\rho^2+g^3/\rho^4)e^{g^2/(2\rho^2)}$.
On a Gaussian slice, the two activation factors are bounded by a
quadratic polynomial in $|b|+\sum_i|w_i|+1$. Bias powers are absorbed
by the bounded Gaussian density times those powers. In the
minimum-coordinate parametrization, integration of polynomial excess
powers is bounded by a constant times the same excess integral,
with the remaining growth at most $(1+t)^5$. Thus the slice majorant
is integrable when
$\int_0^\infty(1+t)^{5-(r-1)}dt<\infty$, which holds for $r\ge8$.
The terms differentiating an activation have at most three powers
and are easier. The bound is uniform in the slice coordinate,
interpolation parameter, and $\beta\ge1$.
Consequently $\beta^2(T_\beta-T)\to D$ by dominated convergence.
Normalize with (31) and subtract
$m_\beta^2\chi\chi^\top$, using (38). This proves (40).
The five-power slice bound fails in dimension six; Theorem~\ref{thm:covariance}
identifies the larger contribution there directly from the critical tail.
\end{proof}

\section*{Dependencies and verification scope}
The argument uses the critical rays and even Walsh block of the preceding
independent parity-posterior papers. It extends the count-and-rank method
of \emph{Polynomial probabilities of extra exceptional neurons} and the
corner calculation of \emph{Smoothing changes the fluctuations of Bayesian
neural condensates}. The uniform temperature estimates and the separate
fine-cutoff argument are the additional steps here. The proofs are
conventional and have not been fully formalized in Lean or independently
reviewed. Conditional one-neuron importance integrations check (19)--(20)
at several logarithmic widths; they do not certify the posterior count
or evidence exponents.
\begin{thebibliography}{9}
\bibitem{cho-saul} Y. Cho and L. K. Saul,
\emph{Kernel methods for deep learning}, Advances in Neural Information
Processing Systems 22 (2009).
\url{https://papers.nips.cc/paper_files/paper/2009/file/5751ec3e9a4feab575962e78e006250d-Paper.pdf}.
\bibitem{sahs} J. Sahs, R. Pyle, F. Anselmi, and A. Patel,
\emph{The spectral bias of shallow neural network learning is shaped
by the choice of non-linearity}, arXiv:2503.10587 (2025).
\url{https://arxiv.org/abs/2503.10587}.
\bibitem{mehrpanah} A. Mehrpanah, M. Gamba, K. Smith, and H. Azizpour,
\emph{On the complexity-faithfulness trade-off of gradient-based explanations},
Proceedings of ICCV (2025), 3531--3541, especially supplement Section F.
\url{https://arxiv.org/abs/2508.10490}.
\end{thebibliography}
\end{document}
